\section{Глава~\ref{sec:recognition}. Распознавание}

Скорость распознавания, устройство первых ступеней, линейное чтение ответа и обучение на непрерывности времени измерены на человеке и обезьяне; сведение всей цепочки к двум операциям и обучение на видео проверены на моделях книги, а объяснения иллюзий --- от хорошо обоснованных до спорных.

{\footnotesize
\setlength{\tabcolsep}{3pt}\renewcommand{\arraystretch}{1.15}
\begin{longtable}{>{\raggedright}p{0.5cm}>{\raggedright}p{4.0cm}>{\raggedright}p{5.6cm}>{\raggedright}p{2.3cm}>{\raggedright\arraybackslash}p{1.7cm}}
\toprule
№ & Утверждение & Опыт и что измерено & Источник & Статус \\
\midrule\endfirsthead
\toprule
№ & Утверждение & Опыт и что измерено & Источник & Статус \\
\midrule\endhead
1 & Сравнение с образцом по пикселям при сдвиге угадывает не лучше монетки; пара ступеней~\eqref{eq:scstage} узнаёт фигуру в любом месте & \texttt{models/\allowbreak recognition\_models.py}, «сетчатка» из $24$ точек, $4000$ проб: образец --- $0.49$, $0.51$, $0.52$ при доле перевёрнутых точек $0$, $0.1$, $0.2$; пара $S$ и $C$ --- $1.00$, $0.86$, $0.70$ & вывод в главе & модель \\
2 & Животное на картинке узнаётся за $\sim150$\,мс, один проход по десятку ступеней по $10$--$20$\,мс & Человек, задача «есть животное или нет» на фотографиях, показанных на $20$\,мс: вызванные потенциалы для двух ответов расходятся через $\sim150$\,мс. Обезьяна: время первого ответа растёт от ступени к ступени (V1 раньше всех, затем V2 и V4). Число ступеней и «ноль или один импульс на ступень» --- вывод из этих чисел & \cite{Thorpe1996,Schmolesky1998} & измерено \\
3 & Ступень $S$ --- И по частям, ступень $C$ --- наибольшее по положениям (утверждение~\ref{prop:conveyor}) & Кошка, первая зрительная кора: простые клетки отвечают на край определённой ориентации в определённом месте, сложные --- на ту же ориентацию в большей области. Внутриклеточная запись сложных клеток: ответ на две полоски у многих клеток близок к большему из двух ответов, в среднем ближе всего к операции max. Наибольшее через взаимное торможение --- утверждение~\ref{prop:wta}, деление на общую активность --- обзор & \cite{HubelWiesel1962,Lampl2004,Carandini2012} & измерено \\
4 & Выше по цепочке сочетания крупнее и сложнее, ответ всё меньше зависит от положения & Обезьяна, упрощение стимулов до «критического признака»: клетки V4 и задней нижневисочной коры отвечают на сложные сочетания формы, цвета и текстуры при малом поле; в передней нижневисочной коре поле больше. Чередование $S$ и $C$ в модели воспроизводит эту картину & \cite{KobatakeTanaka1994,Riesenhuber1999} & измерено \\
5 & Свёрточные сети повторяют это чередование и лучше всех предсказывают ответы верхних ступеней; предмет читается линейно по частотам группы нейронов & Сеть, обученная распознавать, без подгонки к мозгу: верхний слой предсказывает ответы нейронов нижневисочной коры обезьяны, средние --- ответы V4. Линейный классификатор по активности $\sim100$ случайно выбранных клеток за окна от $12.5$\,мс узнаёт предмет и категорию при разном положении и размере & \cite{Fukushima1980,Yamins2014,Hung2005} & измерено \\
6 & Правило Хебба со следом~\eqref{eq:traceinvariance} учит инвариантности на видео без подписей, если след не короче $\sim20$\,мс & Идея медленных признаков и правила со следом --- теория. \texttt{models/\allowbreak recognition\_models.py}, два $C$-нейрона, $16$ входов, $10$ прогонов: пауза между предметами $0.3$\,с. При $\tau_{\text{сл}}=5$\,мс совпадение с фигурой $0.52$ в среднем, при $10$\,мс --- $0.56$; при $20$, $50$ и $200$\,мс --- $1.00$ во всех прогонах (при $30$\,мс один прогон из десяти ошибается) & \cite{Foldiak1991,WiskottSejnowski2002} & модель \\
7 & В мозге инвариантность учится на непрерывности времени & Обезьяна: предмет незаметно подменяли во время скачка глаза в определённое место поля; позиционная инвариантность нейронов нижневисочной коры менялась в соответствии с подменой, значимо уже через час. Что это главное правило обучения зрения --- гипотеза & \cite{LiDiCarlo2008} & измерено \\
8 & Движение: детекторы направления, затем та же пара И/ИЛИ по большим областям & Детекторы направления в сетчатке кролика; в области MT обезьяны все $110$ записанных нейронов избирательны к направлению, ответ на движение сильнее, чем в V1 & \cite{Barlow1965,Albright1984} & измерено \\
9 & Верхние ступени посылают вниз предсказание, вверх идёт ошибка & Модель объясняет внеклассические эффекты полей первой зрительной коры; отдельные наблюдения согласуются (обзор), как общее устройство конвейера не доказано & \cite{RaoBallard1999,Keller2018} & гипотеза \\
10 & Трудные картинки требуют обратных связей и нескольких кругов & Обезьяна: для «трудных» картинок решение в нижневисочной коре складывается на $\sim30$\,мс позже; эти поздние ответы плохо предсказывает прямая сеть и лучше --- сети с обратными связями или более глубокие & \cite{Kar2019} & измерено \\
11 & Незаметная подделка пикселей обманывает сеть; зрение человека гораздо устойчивее, но не неуязвимо & На сетях: малое специально подобранное возмущение меняет ответ; пример «панда $\to$ гиббон» --- из второй работы. У человека при коротком показе подделки, хорошо переносимые между сетями, сдвигают выбор & \cite{Szegedy2014,Goodfellow2015,Elsayed2018} & измерено \\
12 & Иллюзии ожидания: ненадёжный вход уступает ожиданию --- бледное движется «медленнее», платье видят по-разному (раздел~\ref{sec:illusions}) & Решётки меньшего контраста кажутся движущимися медленнее. Опрос $1401$ человека: платье видят в трёх устойчивых вариантах, подсказка об освещении переключает цвет; у $\sim13\,000$ наблюдателей решает предположение об освещении. Байесовская модель с ожиданием медленных движений объясняет ряд иллюзий движения & \cite{Thompson1982,LaferSousa2015,Wallisch2017,Weiss2002,Purves2011} & гипотеза \\
13 & Регулировка усиления: водопад, наклон и контраст; решётка Германа --- не торможение соседей & Детекторы направления сетчатки кролика при долгом движении снижают частоту, а после остановки падают ниже фона. Иллюзия наклона воспроизводится моделью с делением на активность соседей. Изменения решётки, не меняющие ответ выходных нейронов сетчатки, убирают пятна --- объяснение соседним торможением отвергнуто & \cite{BarlowHill1963,Schwartz2009,SchillerCarvey2005} & измерено \\
14 & Компенсация задержек: движущийся предмет кажется впереди вспышки & Иллюзия измерена. Психофизика противоречит и продлению движения вперёд, и разнице задержек; предложено объяснение задним числом --- по событиям $\sim80$\,мс после вспышки. Спор не решён & \cite{Nijhawan1994,EaglemanSejnowski2000} & гипотеза \\
15 & Неподвижные «змеи» вращаются: разница задержек~\eqref{eq:latency} читается детекторами направления & Иллюзия работает и без цвета; пары соседних элементов порождают её поодиночке; нейроны коры обезьяны, избирательные к направлению, дают направленный ответ на эти неподвижные пары. У человека вращение запускают микроскачки глаз и моргания & \cite{Conway2005,OteroMillan2012} & измерено \\
16 & Двойственные картинки: взаимное торможение, усталость и шум; смены вызывает главным образом шум & Модель конкурирующих групп нейронов: смены в ней вызывает \emph{шум}, без шума их нет; она объясняет рост частоты смен с силой стимула. Усталость здесь играет меньшую роль & \cite{MorenoBote2007} & гипотеза \\
17 & Часть иллюзий --- следствие задачи, которой учится машина & Сеть, обученная предсказывать следующий кадр видео, «видит» вращение неподвижных колец и не видит его на контрольных картинках; сети для очистки и повышения резкости изображений поддаются иллюзиям цвета и яркости, как человек & \cite{Watanabe2018,GomezVilla2020} & измерено \\
\bottomrule
\end{longtable}}
